% =====================================================================
% RELATED WORK
% =====================================================================

\section{Related work}
\label{sec:related}

\paragraph{Double descent in deep learning.}
\citet{belkin2019reconciling,nakkiran2021deep} first articulated the
double-descent picture across kernel methods, small neural networks, and
modern deep networks (CNNs, ResNets), showing that the peak in test loss
arises along several axes simultaneously---including model width, training
time, and dataset size---and is amplified by label noise. Subsequent work has refined the picture: sample-wise variants
\citep{nakkiran2020optimal}, the role of regularization in attenuating the
peak \citep{nakkiran2021optimal}, empirical characterizations of when the
phenomenon appears \citep{schaeffer2023double}, and the related grokking
phenomenon \citep{power2022grokking}.
Finally, \citet{wilson2020bayesian} observe empirically that Bayesian
model averaging (approximated by deep ensembles) can eliminate the
double-descent peak. We study this observation by decomposing the MLE
excess test cross-entropy into four terms through the universal-compression
lens---an approximation-and-shift term, a minimax batch regret, a prior cost,
and an MLE-vs-Bayes gap---and, under $100\%$ label noise where the
approximation-and-shift term is zero by construction, find a large
SGD-vs-ensemble gap near interpolation.

\paragraph{Double descent in classical models.}
A complementary line of work analyzes double descent in tractable, mostly
linear, settings: high-dimensional linear regression
\citep{advani2017high,hastie2022surprises}, finite-sample least squares on
concrete data models \citep{belkin2020two}, benign overfitting of
minimum-norm interpolants \citep{bartlett2020benign,muthukumar2020harmless},
ridge regression over random features \citep{mei2022generalization}, and
feature designs exhibiting arbitrarily many descents \citep{chen2021multiple}.
Our contribution is complementary: we target \emph{misspecified, non-regular}
model families, of which deep networks are the prototypical case, rather
than well-specified linear or kernel settings. Unlike these exact risk
analyses, our bound separates candidate contributions without giving a
tight risk formula for deep networks.

\paragraph{Information-theoretic complexity and deep learning theory.}
For discrete data, $-\log q(x)$ is an ideal code length in nats, so expected
excess code length over the true source is a KL divergence.
Classical complexity measures such as VC dimension \citep{vapnik1971uniform}
and Rademacher complexity \citep{bartlett2002rademacher} are notoriously
loose for overparameterized networks; mutual-information generalization
bounds \citep{xu2017information,steinke2020reasoning} provide a
distribution-dependent alternative. Universal coding
\citep{rissanen1984universal,barron1998minimum,davisson1973universal,shtarkov1987universal,haussler1997mutual,merhav1998universal,rissanen1996fisher}
characterizes the unavoidable cost of predicting without knowing the
source parameter, and singular learning theory
\citep{watanabe2009algebraic,hoogland2024towards,chen2023dynamical}
extends this to the non-regular families characteristic of deep networks.
However, none of these frameworks have been shown to directly explain the
double-descent phenomenon in full generality.
